% Source-derived chapter 03: A uniformity estimate for even degree monic polynomials
% Original source: squarefree-values-polynomial-discriminants
\chapter{A uniformity estimate for even degree monic polynomials}
\label{squarefree-values-polynomial-discriminants-chapter-03}
\source{squarefree-values-polynomial-discriminants}

\begin{remark}[Source section]
\label{squarefree-values-polynomial-discriminants-chapter-03-source}
\source{squarefree-values-polynomial-discriminants}
\source{squarefree-values-polynomial-discriminants}
This chapter reproduces the corresponding section of the retrieved
LaTeX source, including its definitions, statements, examples, and
proofs. It has not yet been translated into Lean.
\notready
\end{remark}

% The source section title was: A uniformity estimate for even degree monic polynomials
\label{sec:moniceven}
\source{squarefree-values-polynomial-discriminants}

In this section, which is structured similarly to Section~\ref{sec:monicodd}, we
prove the estimate of Theorem~\ref{thm:mainestimate}(b) when $n=2g+2$
is even, for any $g\geq 1$.

\subsection{Invariant theory for the fundamental representation: $\SO_n$ on the space $W$ of symmetric $n\times n$ matrices}



Let $A_0$ denote the $n\times n$ symmetric matrix with $1$'s on the
anti-diagonal and $0$'s elsewhere. The group $\SO(A_0)$ acts on $W$
via the action
\begin{equation*}
\gamma\cdot B=\gamma B\gamma^t.
\end{equation*}
The central $\mu_2$ acts trivially and so the action descends to an
action of $G=\SO(A_0)/\mu_2$.

We recall some of the arithmetic invariant theory of the
representation $W$ of $n\times n$ symmetric matrices of the
(projective) split orthogonal group $G=\PSO_n.$ See \cite{SW} for more
details. The ring of polynomial invariants over $\C$ is
freely generated by the coefficients of the {\it invariant
  polynomial} $$f_B(x):=(-1)^{g+1}\det(A_0x-B)$$ (See \cite[\S2.1, 2.2]{SW}). We define the {\it
  discriminant} $\Delta$ and {\it height} $H$ on $W$ as the discriminant
and height of the invariant polynomial.

Let $k$ be a field of characteristic not $2$.
%Elements $B\in W(k)$ with $\Delta(B)\neq0$ are stable in the sense of geometric invariant theory.
For any monic polynomial $f(x)\in k[x]$ of degree~$n$ such
that $\Delta(f)\neq0$, let $C_f$ denote the smooth hyperelliptic curve
$y^2=f(x)$ of genus $g$ and let $J_f$ denote its Jacobian. Then $C_f$
has two rational non-Weierstrass points at infinity that are conjugate
by the hyperelliptic involution.
The stabilizer of an element $B\in W(k)$ with invariant polynomial $f(x)$
is isomorphic to $J_f[2](k)$ by~\cite[Proposition~2.33]{W}, and hence has cardinality at most $\#J_f[2](\bar k)=2^{2g}$, where $\bar k$ denotes a separable closure of $k$.

We say that an element (or the $G(k)$-orbit of an element) $B\in W(k)$ with $\Delta(B)\neq 0$ is
{\it distinguished} if there exists a flag $Y'\subset Y$ defined over
$k$ where $Y$ is $(g+1)$-dimensional isotropic with respect to~$A_0$
and $Y'$ is $g$-dimensional isotropic with respect to $B$. If $B$ is
distinguished, then the set of these flags is in bijection with
$J_f[2](k)$ by \cite[Proposition~2.32]{W},  and so it too has cardinality at most~$2^{2g}$.

In fact, it is known (see~\cite[Proposition~22]{BGWhyper}) that the elements of $J_f[2](k)$ are in natural bijection with the even degree factorizations of $f$ defined over $k$.  (Note that the number of such factorizations of $f$ over $\bar k$ is indeed $2^{2g}$.) In particular, if $f$ is irreducible over $k$ and does not factor as $g(x)\bar g(x)$ over some quadratic extension of $k$, then the group $J_f[2](k)$ is trivial.

Let $W_0$ be the subspace of $W$ consisting of matrices whose top left
$g\times (g+1)$ block is zero. Then elements $B$ in $W_0(k)$ with
nonzero discriminant are all distinguished since the $(g+1)$-dimensional
subspace $Y_{g+1}$ spanned by the first $g+1$ basis vectors is
isotropic with respect to $A_0$ and the $g$-dimensional subspace
$Y_g\subset Y_{g+1}$ spanned by the first $g$ basis vectors is
isotropic with respect to $B$.
Let $G_0$ be the parabolic subgroup of
$G$ consisting of elements $\gamma$ such that $\gamma^t$ preserves the
flag $Y_g\subset Y_{g+1}$. Then $G_0$ acts on $W_0$.

An element $\gamma\in G_0$ has the block matrix form
\begin{equation}\label{eq:G_0even}
\source{squarefree-values-polynomial-discriminants}
\gamma=\left(\begin{array}{ccc}\gamma_1 & 0 & 0\\ \delta_1 & \alpha & 0\\ \delta_2 & \delta_3 & \gamma_2
\end{array}\right)\in\left(\begin{array}{ccc}M_{g\times g} & 0 & 0\\ M_{1\times g} & M_{1\times 1} & M_{1\times (g+1)}\\ M_{(g+1)\times g} & M_{(g+1)\times 1} & M_{(g+1)\times (g+1)}
\end{array}\right),
\end{equation}
and so $\gamma\in G_0$ acts on the top right $g\times (g+1)$ block
of an element $B\in W_0$ by
$$\gamma.B^\sub = \gamma_1B^\sub\gamma_2^t,$$ where we use the
superscript ``top'' to denote the top right $g\times (g+1)$ block of
any given element of $W_0$. We may thus view $(A_0^\sub,B^\sub)$ as an
element of the representation $V_g=2\otimes g\otimes (g+1)$. As
before, we define the $Q$-{\it invariant} of $B\in W_0$ as the
$Q$-invariant of $(A_0^\sub,B^\sub)$:
\begin{equation}\label{eqQBeven}
\source{squarefree-values-polynomial-discriminants}
Q(B):=Q(A_0^\sub,B^\sub).
\end{equation}
Then the $Q$-invariant is a relative invariant for the action of $G_0$ on $W_0$, i.e.,
for any $\gamma\in G_0$ in the form \eqref{eq:G_0even}, we have by \eqref{eq:relQ},
\begin{equation}\label{eq:weightG_0even}
\source{squarefree-values-polynomial-discriminants}
Q(\gamma.B) = \det(\gamma_1)^{g+1}\det(\gamma_2)^gQ(B) =
\det(\gamma_1)\alpha^{-g}Q(B).
\end{equation}

In fact, we may extend the definition of the
$Q$-invariant to an even larger subset of $W(\Q)$ than $W_0(\Q)$.
We have the following proposition.

\begin{proposition}
\source{squarefree-values-polynomial-discriminants}
\notready\label{prop:extendQeven}
\source{squarefree-values-polynomial-discriminants}
Let $B\in W_0(\Q)$ be an element whose invariant polynomial $f(x)$ is
irreducible over $\Q$ and, when $n\geq 4$, does not factor as $g(x)\bar{g}(x)$ over some quadratic extension of $\Q$. Then for every $B'\in W_0(\Q)$ such that $B'$ is $G(\Z)$-equivalent to $B$, we have
$Q(B')=\pm Q(B)$.
\end{proposition}
\begin{proof}
The assumption on the factorization property of $f(x)$ implies that $J_f[2](\Q)$ is trivial.
The proof is now identical to that of Proposition~\ref{prop:extendQ}.
\end{proof}

We may thus define the $|Q|$-{\it invariant} for any element $B\in W(\Q)$ that is $G(\Z)$-equivalent to some $B'\in W_0(\Q)$ and whose invariant polynomial is irreducible over $\Q$ and does not factor as $g(x)\bar{g}(x)$ over any quadratic extension of $\Q$; indeed, we set $|Q|(B):=|Q(B')|$. By Proposition~\ref{prop:extendQeven}, this definition of $|Q|(B)$ is independent of  the choice of $B'$. We note again that all such elements $B\in W(\Q)$ are distinguished.

\subsection{Embedding $\U_m^\w$ into $\frac14W(\Z)$}\label{sembedeven}
\source{squarefree-values-polynomial-discriminants}

Let $m$ be a positive squarefree integer and let $f$ be an monic
integer polynomial whose discriminant is weakly divisible by
$m^2$. Then as proved in \S\ref{sembedodd}, there exists an integer
$\ell$ such that $f(x+\ell)$ has the form
$$f(x+\ell) = x^n + c_1x^{n-1} + \cdots + c_{n-2}x^2 + mc_{n-1}x + m^2c_n.$$



Consider the following matrix:
\begin{equation}\label{mat2}
\source{squarefree-values-polynomial-discriminants}
 B_m(c_1,\ldots,c_n) =
 \left(\!\!\!\!\!\begin{array}{cccccccccc}&&&&&&&&m&0\\[.065in]&&&&&&&\iddots&\:\;\iddots& \\[.025in]&&&&&&1&\;\;\iddots&&\\[.185in] &&&&&\,1&0&&&\\[.185in] &&&&\:\;1\;&-c_1/2&&&& \\[.175in]&&&1&\!\!-c_1/2\,&\!\!c_1^2/4\!-\!c_2\!\!&-c_3/2&&&\\[.175in]&&1&\;\;\;0\;\;\;&&-c_3/2&-c_4&\!\!\!-c_5/2\!\!\!&&\\[.085in]&\iddots&\;\;\:\iddots\:\;\;&&&&-c_5/2&-c_6&\ddots&\\[.0125in] \;\;\;m\;\;\:&\;\,\,\iddots&&&&&&\ddots&\ddots&\!\!-c_{n-1}/2\!\!
   \\[.125in] 0&&&&&&&&\!\!\!\!-c_{n-1}/2\!\!\!\!&-c_n \end{array}\!\right).
\end{equation}
It follows from a direct computation that
$$f_{B_m(c_1,\ldots,c_n)}(x) = x^n + c_1x^{n-1} + \cdots + c_{n-2}x^2 +
mc_{n-1}x + m^2c_n.$$ We set $\sigma_m(f) := B_m(c_1,\ldots,c_n) +
\ell A_0\in \frac14 W(\Z)$. Then evidently $f_{\sigma_m(f)}=f.$ A direct computation again shows that
$Q(B_m(c_1,\ldots,c_n))=m$. Since the $Q$-invariant on $2\otimes
g\otimes (g+1)$ is $\SL_2$-invariant, we conclude
that $$Q(\sigma_m(f))=m.$$


\begin{theorem}
\source{squarefree-values-polynomial-discriminants}
\notready\label{th:mapeven}
\source{squarefree-values-polynomial-discriminants}
Let $m$ be a positive squarefree integer. There exists a map
$\sigma_m:\U_m^\w\to \frac14W(\Z)$ such that for every $f\in \U_m^\w$,
\begin{equation}\label{eq:liftQe}
\source{squarefree-values-polynomial-discriminants}
f_{\sigma_m(f)}=f,\qquad Q(\sigma_m(f)) = m.
\end{equation}
\end{theorem}

Let $\LL$ be the set of elements $v\in \frac14W(\Z)$ that are $G(\Z)$-equivalent to some elements of $\frac14 W_0(\Z)$ and such that
the invariant polynomial of $v$ is irreducible
over $\Q$ and does not factor as $g(x)\bar{g}(x)$ over some quadratic extension of $\Q$.
Then by the remark following Proposition~\ref{prop:extendQeven}, we may view $|Q|$ as a function also on $\LL$.
Let $\U_m^{\w,\irr}$ denote the set of polynomials in $\U_m^\w$
that are irreducible
over $\Q$ and do not factor as $g(x)\bar{g}(x)$ over any quadratic extension of $\Q$. Then we have the following immediate consequence of Theorem~\ref{th:mapeven}:

\begin{theorem}
\source{squarefree-values-polynomial-discriminants}
\notready\label{keymaporbiteven}
\source{squarefree-values-polynomial-discriminants}
Let $m$ be a positive squarefree integer. There exists an injective map
$$\bar{\sigma}_m:\U_m^{\w,\irr}\to G(\Z)\backslash\LL$$ such that
$f_{\bar{\sigma}_m(f)}=f$ for every $f\in \U_m^\w$. Furthermore, every element
in every orbit in the image of $\bar{\sigma}_m$ has $|Q|$-invariant $m$.
\end{theorem}

The number of monic integer polynomials having height less than $X$ that are reducible or factor as $g(x)\bar{g}(x)$ over some quadratic extension of $\Q$ is of a strictly smaller order of magnitude than the total number of such polynomials (see, e.g., Proposition~\ref{propredboundall}). Thus to prove Theorem~\ref{thm:mainestimate}(b), it suffices to count the number of elements in
$\U_m^{\w,\irr}$ having height less than $X$ over all $m>M$, which, by Theorem~\ref{keymaporbiteven}, we may do by counting
$G(\Z)$-orbits on $\LL\subset \frac14W(\Z)$ having height less than $X$ and $|Q|$-invariant greater than $M$.
More precisely, let $N(\LL;M;X)$ denote the number of $G(\Z)$-equivalence classes of elements in
$\LL$ whose $|Q|$-invariant is greater than $M$ and whose height is
less than $X$. We obtain a bound for $N(\LL;M;X)$ using the same method as in Section 2.



\subsection{Counting $G(\Z)$-orbits in $\frac14W(\Z)$}\label{scoeffeven}
\source{squarefree-values-polynomial-discriminants}

The counting problem for the representation $W$ of $G$ is studied in \cite{SW}. In this section, we recall some of the set up and results of \cite{SW}.

Let $R$ be a fundamental domain for the action of $G(\R)$ on the elements of $W(\R)$ having nonzero discriminant and height bounded by $1$ as constructed in \cite[\S4.1]{SW}. Let $\FF$ be a fundamental set for the left multiplication action of $G(\Z)$ on $G(\R)$ obtained using the Iwasawa decomposition of $G(\R)$. More explicitly, we have
$$
G(\R)=N(\R)TK,
$$ where $N$ is a unipotent group consisting of lower triangular matrices, $K$ is compact, and $T$ is the split torus of $G$ given by
\begin{equation*}
T=
\left\{\left(\begin{array}{ccccccc}
 t_1^{-1}&&&&&&\\
&\ddots &&&&&  \\
 && t_{g+1}^{-1} &&&&\\
 &&&& \!\!\!\!t_{g+1} &&\\
&&&&&\!\!\!\!\ddots &  \\
& &&&&& \!\!t_{1}
\end{array}\right)\right\}.
\end{equation*}
We may also make the following change of variables.
For $1\leq
i\leq g$, define $s_i$ to be
$$
s_i=t_i/t_{i+1},
$$
and let $s_g=t_gt_{g+1}$.
We denote an element of $T$ with coordinates $t_i$
(resp.\ $s_i$) by $(t)$ (resp.\ $(s)$). We may take $\FF$ to be contained in a {\it Siegel set}, i.e., contained in $N'T'K$,
where $N'$ consists of elements in $N(\R)$ whose coefficients are
absolutely bounded and $T'\subset T$ consists of elements in $(s)\in
T$ with $s_i\geq c$ for some positive constant $c$.

Let $\GG_1$ be a compact left $K$-invariant set in $G(\R)$ which is the closure of nonempty open set. Then as in Section 2, we have
\begin{equation}
  N(\LL;M;X)\ll \int_{\gamma\in\FF} \#\{B\in((\gamma \GG_1)\cdot (XR))\cap\LL:|Q|(B)>M\}
  d\gamma,
\end{equation}
where the implied constant depends only on $\GG_1$ and $R$, and where $d\gamma$ is a Haar measure on $G(\R)$ given by
$$
d\gamma=dn\,\delta(s)d^\times s\,dk,
$$ where $dn$ is a Haar measure on the unipotent group $N(\R)$,
$dk$ is a Haar measure on the compact group~$K$, $d^\times s$ is
given by
$$
d^\times s:=\prod_{i=1}^{g+1}\frac{ds_i}{s_i},
$$
and
$\delta(s)$ is given by
\begin{equation}\label{eqhaareven}
\source{squarefree-values-polynomial-discriminants}
  \delta(s)=\prod_{k=1}^{g-1} s_k^{k^2-2kg-k}\cdot (s_gs_{g+1})^{-g(g+1)/2};
\end{equation}
see~\cite[(20)]{SW}.


Since $s_i\geq c$ for every $i$, there exists a compact subset $N''$ of
$N(\R)$ containing $(t)^{-1}N'\,(t)$ for all $t\in T'$. Since $N''$, $K$, $G_0$ are compact and $R$ is bounded, the set $E=N''KG_0R$ is bounded. Then we have
\begin{equation}\label{eq:REe}
\source{squarefree-values-polynomial-discriminants}
N(\LL;M;X)\ll \int_{s_i\gg 1} \#\{B\in((s)\cdot XE)\cap\LL:|Q|(B)>M\}\delta(s)d^\times s.
\end{equation}

As before, we denote the coordinates of $W$ by $b_{ij}$, for $1\leq
i\leq j\leq n$, and we denote
the $T$-weight of a coordinate $\alpha$ on $W$, or a
product $\alpha$ of powers of such coordinates, by $w(\alpha)$.
We compute the weights of the coefficients $b_{ij}$ to be
\begin{equation}\label{wbij2}
\source{squarefree-values-polynomial-discriminants}
w(b_{ij})=\left\{
\begin{array}{rcl}
t_i^{-1}t_j^{-1} &\mbox{ if }& i,j\leq g+1\\
t_i^{-1}t_{n-j+1} &\mbox{ if }& i\leq g+1,\; j>g+1\\
t_{n-i+1}t_{n-j+1} &\mbox{ if }& i,j>g+1.
\end{array}
\right.
\end{equation}
Then the $(i,j)$-entry of any $B\in (s)\cdot XE$ is bounded by $Xw(b_{ij})$, up to a multiplicative constant depending only on $\GG_1$ and $R$.

Let $W_{00}\subset W$ denote the space of symmetric matrices $B$
such that $b_{ij}=0$ for $i+j < n$. Let $W^\dist$ denote the subset of $W(\Q)$ consisting of $\Q$-distinguished elements. Then $\LL$ is a subset of $W^\dist$. It was shown in \cite[\S4.2]{SW}
that most of the lattice points in $W^\dist$ lie inside $W_{00}$. More precisely, we have the following estimates from \cite[Proposition 21 and 23]{SW}.

\begin{proposition}
\source{squarefree-values-polynomial-discriminants}
\notready\label{prop:sw6}
\source{squarefree-values-polynomial-discriminants}
We have
\begin{eqnarray}
% \nonumber % Remove numbering (before each equation)
  \int_{s_i\gg 1} \#\{B\in((s)\cdot XE)\cap({\textstyle\frac14} W(\Z)\setminus{\textstyle\frac14} W_{00}(\Z)):b_{11}=0\}\delta(s)d^\times s &=& O_\epsilon(X^{n(n+1)/2-1+\epsilon}) \\
  \label{eq:selbergpree}\int_{s_i\gg 1} \#\{B\in((s)\cdot XE)\cap\LL:b_{11}\neq 0\}\delta(s)d^\times s &=& o(X^{n(n+1)/2}).
\source{squarefree-values-polynomial-discriminants}
\end{eqnarray}
\end{proposition}

Therefore, to prove Theorem~\ref{thm:mainestimate}(b) when $n=2g+2$ is even, it remains to obtain a power saving improvement of \eqref{eq:selbergpree} and estimate
\begin{equation}\label{eq:Qinve}
\source{squarefree-values-polynomial-discriminants}
\int_{s_i\gg 1} \#\{B\in ((s)\cdot XE)\cap \LL\cap
{\textstyle\frac14} W_{00}(\Z):|Q|(B)>M\}\delta(s)d^\times s.
\end{equation}

\subsection{Proof of Theorem~\ref{thm:mainestimate}(b) for even $n$}\label{sgomeven}
\source{squarefree-values-polynomial-discriminants}

We begin with a power saving improvement of \eqref{eq:selbergpree}.

\begin{proposition}
\source{squarefree-values-polynomial-discriminants}
\notready\label{propeven}
\source{squarefree-values-polynomial-discriminants}
  We have
  \begin{equation*}
\displaystyle\int_{s_i\gg 1} \#\{B\in ((s)\cdot XE)\cap W^\dist \cap {\textstyle\frac14} W(\Z): b_{11}\neq 0\}\delta(s)d^\times s =O_\epsilon(X^{n(n+1)/2-1/5+\epsilon}).
  \end{equation*}
\end{proposition}
\begin{proof}
In the proof of \cite[Proposition 23]{SW}, it is shown that the
$p$-adic density of elements in $W(\Z_p)$ that are
$\Q_p$-distinguished is bounded uniformly away from $1$.  Then an
application of the Selberg sieve exactly as in \cite{ShTs} yields the
result.
%We note by \cite[Proposition 23]{SW} that the $p$-adic density of
%elements in $W(\Z_p)$ that are $\Q_p$-distinguished is bounded above
%by a constant $c_n$, depending only on $n$, strictly less than 1. Then
%an application of the Selberg sieve exactly as in \cite{ShTs} yields
%the result.
\end{proof}

We now estimate \eqref{eq:Qinve}. The $Q$-invariant can also be given
a weight by viewing the torus $T$ as sitting inside $G_0$ and using
\eqref{eq:weightG_0even}. Namely,
\begin{equation}\label{eqQweighteven}
\source{squarefree-values-polynomial-discriminants}
w(Q)=(t_1\cdots t_g)^{-1}t_{g+1}^g=\prod_{k=1}^g s_k^k.
\end{equation}
Since the polynomial $Q$ is homogeneous of degree $g(g+1)/2$ in the
coefficients of $W_0$, we see that the $Q$-invariant of any $B\in
(s)\cdot XE$ is bounded by $X^{g(g+1)/2}w(Q)$, up to a multiplicative
constant depending only on $G_0$ and $R$.


\begin{proposition}
\source{squarefree-values-polynomial-discriminants}
\notready\label{proplargeQboundeven}
\source{squarefree-values-polynomial-discriminants}
We have
\begin{equation*}
\int_{s_i\gg 1} \#\{B\in ((s)\cdot(XE))\cap \LL\cap {\textstyle\frac14} W_{00}(\Z):|Q|(B)>M\}\delta(s)d^\times s=O(X^{n(n+1)/2}\log^2 X/M).
\end{equation*}
\end{proposition}
\begin{proof}
Analogous to the proof of
Proposition \ref{proplargeQbound}, in order for the set
$\{B\in((s)\cdot (XE))\cap \LL\cap \frac14 W_{00}(\Z):|Q|(B)>M\}$ to be nonempty,
the following conditions must be satisfied:
\begin{equation}\label{eqscond}
\source{squarefree-values-polynomial-discriminants}
\begin{array}{rcl}
Xs_i^{-1}&\gg& 1,\\[.1in]
Xs_gs_{g+1}^{-1}&\gg& 1,\\[.1in]
X^{g(g+1)/2}w(Q)&\gg& M.
\end{array}
\end{equation}

Let $S$ denote the set of coordinates of $W_{00}$, i.e.,
$S=\{b_{ij}:i+j\geq n\}$. Let $T_{X,M}$ denote the set of~$(s)$
satisfying $s_i\gg 1$ and the conditions of \eqref{eqscond}. Then we have
\begin{eqnarray*}
&&\displaystyle\int_{s_i\gg 1} \#\{B\in((s)\cdot (XE))\cap \LL\cap {\textstyle\frac12} W_{00}(\Z):|Q|(B)>M\}\delta(s)d^\times s \\
&\ll&\displaystyle\int_{(s)\in T_{X,M}} \big(\prod_{\alpha\in S}(Xw(\alpha))\big)\delta(s)d^\times s \\
&\ll&\displaystyle\int_{(s)\in T_{X,M}} X^{n(n+1)/2-g(g+1)}\prod_{k=1}^{g-1}s_k^{2k-1}\cdot s_g^{g-1}s_{g+1}^{g}d^\times s \\
&\ll&\displaystyle\frac{1}{M}\int_{(s)\in T_{X,M}} X^{n(n+1)/2-g(g+1)/2}w(Q)\prod_{k=1}^{g-1}s_k^{2k-1}\cdot s_g^{g-1}s_{g+1}^{g}d^\times s \\
&\ll&\displaystyle\frac{1}{M}\int_{(s)\in T_{X,M}} X^{n(n+1)/2-g(g+1)/2}\prod_{k=1}^{g-1}s_k^{k-1}\cdot s_g^{-1}s_{g+1}^{g}d^\times s \\
&\ll&\displaystyle\frac{1}{M}\int_{(s)\in T_{X,M}} X^{n(n+1)/2-g(g+1)/2+g}\prod_{k=1}^{g-1}s_k^{k-1}\cdot s_g^{g-1}d^\times s \\
&\ll&\displaystyle\frac{1}{M}X^{n(n+1)/2}\log^2(X),
\end{eqnarray*}
where the first inequality follows from the fact that $Xw(b_{ij})\gg1$ for all $b_{ij}\in S$ when $(s)$ is in the range $1\ll s_i\ll X$, the second inequality follows from the definition \eqref{eqhaareven} of $\delta(s)$ and the computation (\ref{wbij2}) of the weights of the coordinates $b_{ij}$, the third inequality follows from the fact that $X^{g(g+1)/2}w(Q) \gg M$,
the fourth inequality follows from the computation of the weight
of $Q$ in \eqref{eqQweighteven}, the fifth inequality comes from multiplying by the factor $(Xs_gs_{g+1}^{-1})^g\gg1$, and  the $\log^2
X$ factor in the last inequality comes from the integrals over $s_1$ and $s_{g+1}$.
\end{proof}


The estimate in Theorem \ref{thm:mainestimate}(b) for even $n$ now
follows from Theorem \ref{keymaporbiteven} and Propositions~\ref{prop:sw6},
\ref{propeven} and \ref{proplargeQboundeven}, in conjunction with the
bound on the number of reducible polynomials proved in
Proposition~\ref{propredboundall}.
